\documentclass[a4paper,11pt]{article}

% ─── Packages ────────────────────────────────────────────────────────────────
\usepackage[left=1.1in,right=1.1in,top=1in,bottom=1in]{geometry}
\usepackage{amsmath}
\usepackage{amssymb}
\usepackage[utf8]{inputenc}
\usepackage[T1]{fontenc}
\usepackage{lmodern}
\usepackage{microtype}

% Colors
\usepackage{xcolor}
\definecolor{primaryblue}{HTML}{1E3A5F}
\definecolor{accentblue}{HTML}{2E86C1}
\definecolor{lightgray}{HTML}{F4F6F9}
\definecolor{medgray}{HTML}{BDC3C7}
\definecolor{darkgray}{HTML}{2C3E50}
\definecolor{codebg}{HTML}{F0F3F7}
\definecolor{codeframe}{HTML}{C5D3E0}
\definecolor{noteframe}{HTML}{D6EAF8}
\definecolor{notebg}{HTML}{EBF5FB}
\definecolor{kw}{HTML}{1A5276}
\definecolor{cm}{HTML}{1E8449}
\definecolor{st}{HTML}{922B21}
\definecolor{ln}{HTML}{95A5A6}

% Links
\usepackage{hyperref}
\hypersetup{
    colorlinks=true,
    linkcolor=accentblue,
    urlcolor=accentblue,
    citecolor=accentblue,
    pdftitle={RepoExplorer --- Technical Documentation},
    pdfauthor={Sahil Arun Dhawane}
}

% Graphics & Tables
\usepackage{graphicx}
\usepackage{booktabs}
\usepackage{array}
\usepackage{tabularx}

% Lists
\usepackage{enumitem}
\setlist[itemize]{leftmargin=1.4em, itemsep=2pt, topsep=4pt}
\setlist[enumerate]{leftmargin=1.6em, itemsep=2pt, topsep=4pt}

% Code listings
\usepackage{listings}
\lstset{
    basicstyle=\ttfamily\footnotesize,
    breaklines=true,
    breakatwhitespace=false,
    frame=single,
    framesep=6pt,
    rulecolor=\color{codeframe},
    backgroundcolor=\color{codebg},
    keywordstyle=\bfseries\color{kw},
    commentstyle=\itshape\color{cm},
    stringstyle=\color{st},
    numbers=left,
    numberstyle=\tiny\color{ln},
    numbersep=8pt,
    showstringspaces=false,
    tabsize=2,
    xleftmargin=1.5em,
    xrightmargin=0.5em,
    aboveskip=8pt,
    belowskip=8pt
}

% Framed boxes
\usepackage{mdframed}
\mdfdefinestyle{notebox}{
    linecolor=accentblue,
    linewidth=1.2pt,
    backgroundcolor=notebg,
    innertopmargin=8pt,
    innerbottommargin=8pt,
    innerleftmargin=12pt,
    innerrightmargin=12pt,
    skipabove=10pt,
    skipbelow=8pt,
    roundcorner=4pt
}
\mdfdefinestyle{flowbox}{
    linecolor=medgray,
    linewidth=0.8pt,
    backgroundcolor=lightgray,
    innertopmargin=8pt,
    innerbottommargin=8pt,
    innerleftmargin=12pt,
    innerrightmargin=12pt,
    skipabove=10pt,
    skipbelow=8pt,
    roundcorner=3pt
}

% Section formatting
\usepackage{titlesec}
\titleformat{\section}
    {\Large\bfseries\color{primaryblue}}
    {}
    {0em}
    {}
    [\vspace{2pt}\textcolor{accentblue}{\rule{\linewidth}{1.5pt}}]
\titlespacing{\section}{0pt}{22pt}{10pt}

\titleformat{\subsection}
    {\large\bfseries\color{darkgray}}
    {}
    {0em}
    {}
\titlespacing{\subsection}{0pt}{14pt}{5pt}

\titleformat{\subsubsection}
    {\normalsize\bfseries\color{darkgray}}
    {}
    {0em}
    {}
\titlespacing{\subsubsection}{0pt}{10pt}{3pt}

% Header / Footer
\usepackage{fancyhdr}
\pagestyle{fancy}
\fancyhf{}
\renewcommand{\headrulewidth}{0.6pt}
\renewcommand{\footrulewidth}{0.4pt}
\fancyhead[L]{\small\textcolor{darkgray}{\textbf{Sahil Arun Dhawane}}}
\fancyhead[R]{\small\textcolor{darkgray}{RepoExplorer --- Technical Documentation}}
\fancyfoot[R]{\small\textcolor{medgray}{Page \thepage}}
\fancyfoot[L]{\small\textcolor{medgray}{github.com/Sahil0282/repo-explorer}}

% TOC styling
\usepackage{tocloft}
\renewcommand{\cftsecfont}{\bfseries\color{primaryblue}}
\renewcommand{\cftsecpagefont}{\bfseries\color{primaryblue}}
\renewcommand{\cftsubsecfont}{\color{darkgray}}
\setlength{\cftbeforesecskip}{4pt}
\setlength{\cftbeforesubsecskip}{2pt}

% Paragraph spacing
\setlength{\parskip}{5pt}
\setlength{\parindent}{0pt}

% ─── Document ─────────────────────────────────────────────────────────────────
\begin{document}

% ═══════════════════════════════════════════════════════
%  TITLE PAGE
% ═══════════════════════════════════════════════════════
\begin{titlepage}
    \centering
    \vspace*{2.5cm}

    {\fontsize{42}{50}\selectfont\bfseries\color{primaryblue} RepoExplorer}\\[0.6cm]
    \textcolor{accentblue}{\rule{0.6\linewidth}{1.5pt}}\\[0.8cm]

    {\Large\color{darkgray} AI-Powered Codebase Understanding Tool}\\[0.4cm]
    {\normalsize\color{medgray} Full Technical Documentation}\\[2.5cm]

    {\large\bfseries\color{darkgray} Sahil Arun Dhawane}\\[0.25cm]
    {\normalsize Vishwakarma Institute of Technology, Pune}\\[0.15cm]
    {\normalsize B.Tech Electronics \& Telecommunications}\\[2.5cm]

    \textcolor{medgray}{\rule{0.4\linewidth}{0.4pt}}\\[0.5cm]
    {\small \textbf{Repository:} \href{https://github.com/Sahil0282/repo-explorer}{github.com/Sahil0282/repo-explorer}}\\[1cm]
    {\small\itshape Prepared for placement interview preparation}\\[0.4cm]
    {\small April 2026}

    \vfill
\end{titlepage}

% ═══════════════════════════════════════════════════════
%  TABLE OF CONTENTS
% ═══════════════════════════════════════════════════════
\tableofcontents
\newpage

% ═══════════════════════════════════════════════════════
%  SECTION 1: PROBLEM STATEMENT
% ═══════════════════════════════════════════════════════
\section{The Problem: Why This Tool Exists}

\subsection{The Core Pain Point}

Every software developer has faced this situation: you are assigned to fix a bug or add a feature in a codebase you have never seen before. Or you want to use an open-source library and need to understand how it works internally. Or you join a new team and are handed a repository with hundreds of files and zero documentation.

In all of these cases, the process of understanding the codebase is brutally slow and frustrating. You open the repository on GitHub and see this:

\begin{mdframed}[style=notebox]
\textbf{Example Scenario:} You need to understand how authentication works in a Node.js backend project. The repository has 200 files across 15 folders. You have no idea where to start. Is it in \texttt{middleware/auth.js}? Or \texttt{utils/jwt.js}? Or \texttt{controllers/userController.js}? Or is it split across all three? You have to manually open each file, read through hundreds of lines of code, and slowly build a mental map of how everything connects. This process can take hours or even days.
\end{mdframed}

\subsection{Why Existing Solutions Fall Short}

\subsubsection{GitHub's Built-in Navigation}
GitHub lets you browse files and view code. But it only shows you \textit{what} the code looks like. It cannot tell you \textit{how} it works, \textit{why} it's structured this way, or \textit{which files you should look at} to understand a specific feature. You still need to already know which file to open.

\subsubsection{IDE Tools like Copilot and Cursor}
Tools like GitHub Copilot Chat and Cursor are excellent for developers who are \textit{already working inside} a codebase. They require you to:
\begin{itemize}
    \item Clone the repository to your local machine
    \item Open it in a supported IDE (VS Code, JetBrains, etc.)
    \item Wait for indexing to complete
    \item Have the right extensions installed
\end{itemize}
This is a 15--30 minute setup process minimum. And it only works if you \textit{own} the codebase or have local access to it.

\subsubsection{Documentation}
Most codebases have little to no documentation. Even when documentation exists, it is often outdated and does not reflect the current state of the code.

\subsection{The Gap RepoExplorer Fills}

RepoExplorer targets a fundamentally different user moment:

\begin{mdframed}[style=notebox]
\textbf{The user does not own the code. They are outside the codebase, looking in.}
\end{mdframed}

\noindent Examples of this user:
\begin{itemize}
    \item A new joiner who has not been given laptop access yet
    \item An open-source contributor evaluating whether a library suits their needs
    \item A tech lead reviewing a candidate's GitHub project during an interview
    \item A developer asked to debug someone else's project without any context
    \item A student trying to understand how a popular open-source project is structured before contributing
\end{itemize}

For all of these people, the existing tools are useless. RepoExplorer solves this by allowing any user to paste a GitHub URL, ask a question in plain English, and receive a structured answer with exact file citations and a visual map of which functions are involved --- all in under 60 seconds with zero local setup.

\subsection{Concrete Example of the Problem and Solution}

\textbf{Without RepoExplorer:}\\
User wants to understand how authentication works in \texttt{expressjs/express}. They open GitHub, browse to the repository, see 213 files, start opening random files, spend 45 minutes reading code, still are not sure which middleware runs first.

\vspace{6pt}
\textbf{With RepoExplorer:}\\
User pastes \texttt{https://github.com/expressjs/express} into the tool. Asks: \textit{``how does authentication work in this project?''}. Within 5 seconds receives:

\begin{mdframed}[style=notebox]
Authentication in this project involves the following flow:\\[4pt]
1.\ \textbf{examples/auth/index.js (lines 25--34):} The \texttt{loadUser} middleware fetches a user based on \texttt{req.params.id} and attaches it to \texttt{req.user}\ldots\\[2pt]
2.\ \textbf{controller/authController.js (lines 56--64):} Upon successful sign-up, an \texttt{authtoken} is set as an HTTP cookie and included in the JSON response\ldots\\[2pt]
\ldots and so on with exact citations.
\end{mdframed}

The user also sees an Execution Flow on the right side showing exactly how the relevant functions connect, and can click any function node to jump to its exact code.

\newpage

% ═══════════════════════════════════════════════════════
%  SECTION 2: SYSTEM OVERVIEW
% ═══════════════════════════════════════════════════════
\section{System Overview: How RepoExplorer Solves This}

\subsection{High-Level Architecture}

RepoExplorer is built as three separate services that work together:

\begin{enumerate}
    \item \textbf{Frontend} --- React application (runs on port 5173)
    \item \textbf{Node.js Backend} --- Express server (runs on port 8000)
    \item \textbf{Python AI Service} --- FastAPI server (runs on port 8001)
\end{enumerate}

\noindent\textbf{Flow when a user analyzes a repository:}

\begin{mdframed}[style=flowbox]
\ttfamily\small
User (Browser) $\rightarrow$ React Frontend $\rightarrow$ POST /api/repo/analyze (Node.js:8000)\\
\quad$\rightarrow$ Clone repo from GitHub (simple-git)\\
\quad$\rightarrow$ Parse all JS/TS files (Babel AST)\\
\quad$\rightarrow$ Extract all functions with code + metadata\\
\quad$\rightarrow$ Compute git evolution map (gitEvolution.js) \textbf{[NEW]}\\
\quad$\rightarrow$ POST /index (Python FastAPI:8001)\\
\quad$\rightarrow$ Build enriched chunks\\
\quad$\rightarrow$ Generate embeddings (MiniLM-L6-v2)\\
\quad$\rightarrow$ Store in ChromaDB\\
\quad$\rightarrow$ Return file tree + stats to Frontend
\end{mdframed}

\noindent\textbf{Flow when a user asks a question:}

\begin{mdframed}[style=flowbox]
\ttfamily\small
User types question $\rightarrow$ POST /api/repo/query (Node.js:8000)\\
\quad$\rightarrow$ POST /query (Python FastAPI:8001)\\
\quad$\rightarrow$ Hybrid Search: BM25 + Vector + RRF\\
\quad$\rightarrow$ Top 5--8 relevant chunks retrieved\\
\quad$\rightarrow$ Chunks + question sent to Gemini 2.5 Flash\\
\quad$\rightarrow$ Structured JSON returned: markdown answer + execution\_flow graph \textbf{[UPDATED]}\\
\quad$\rightarrow$ Answer + source files + source functions + execution flow sent to Frontend\\
\quad$\rightarrow$ Chat message rendered with markdown (file/function citations clickable) \textbf{[UPDATED]}\\
\quad$\rightarrow$ Execution Flow graph rendered on demand (React Flow + Dagre) \textbf{[UPDATED]}
\end{mdframed}

\noindent\textbf{Flow when a user requests function evolution history:} \textbf{[NEW]}

\begin{mdframed}[style=flowbox]
\ttfamily\small
User requests evolution $\rightarrow$ POST /api/repo/evolution (Node.js:8000)\\
\quad$\rightarrow$ gitEvolution.js reads last 20 commits (git show {-}{-}unified=0)\\
\quad$\rightarrow$ Parse added/deleted lines per commit with line numbers\\
\quad$\rightarrow$ Map line changes $\rightarrow$ extracted functions by startLine/endLine range\\
\quad$\rightarrow$ Build evolutionMap keyed as filePath::functionName\\
\quad$\rightarrow$ POST /evolution-summary (Python FastAPI:8001)\\
\quad$\rightarrow$ Gemini summarizes change history per function\\
\quad$\rightarrow$ Structured evolution timeline returned to Frontend
\end{mdframed}

\begin{mdframed}[style=notebox]
\textbf{Important distinction:} Git commit diffs are \textbf{not embedded}. They are sent directly to Gemini as text for summarization. Only code chunks (function bodies + file summaries) go through MiniLM embeddings and ChromaDB. The two pipelines are completely independent.
\end{mdframed}

\newpage

% ═══════════════════════════════════════════════════════
%  SECTION 3: AI SERVICE IN DETAIL
% ═══════════════════════════════════════════════════════
\section{The AI Service: Every Component Explained}

This is the heart of the system. The Python FastAPI service handles all AI-related operations. It lives in the \texttt{service/} folder.

\subsection{The RAG (Retrieval-Augmented Generation) Paradigm}

\subsubsection{What is RAG and Why We Use It}

Large Language Models like Gemini are trained on massive datasets but have two fundamental limitations:
\begin{enumerate}
    \item \textbf{Knowledge cutoff:} They don't know about code written after their training cutoff.
    \item \textbf{Context window:} You cannot feed an entire codebase (potentially millions of lines) into an LLM prompt --- it has a token limit.
\end{enumerate}

RAG solves both problems. Instead of sending the entire codebase to the LLM, you:
\begin{enumerate}
    \item \textbf{Index phase:} Break the codebase into small chunks, convert each chunk to a vector (embedding), store in a vector database.
    \item \textbf{Query phase:} When a question arrives, find the most relevant chunks using similarity search, send only those chunks to the LLM.
\end{enumerate}

The LLM then answers based on the retrieved chunks --- it is grounded in real code, not hallucinating. This is the core principle behind RepoExplorer.

\subsection{Embedding Model: MiniLM-L6-v2}

\subsubsection{What is an Embedding?}
An embedding is a numerical representation of text as a vector (list of floating point numbers). Text with similar meaning gets vectors that are mathematically close to each other in high-dimensional space.

Example:
\begin{itemize}
    \item \texttt{"function that handles user authentication"} $\rightarrow$ [0.23,\;$-$0.45,\;0.12,\;\ldots]
    \item \texttt{"JWT token verification middleware"} $\rightarrow$ [0.21,\;$-$0.43,\;0.15,\;\ldots]
    \item \texttt{"sorting a list of integers"} $\rightarrow$ [$-$0.67,\;0.89,\;$-$0.34,\;\ldots]
\end{itemize}

The first two vectors are close together (similar meaning), the third is far away.

\subsubsection{Why MiniLM-L6-v2 Specifically?}
We use the \texttt{sentence-transformers/all-MiniLM-L6-v2} model from Hugging Face. Here is why this model was chosen over alternatives:

\vspace{6pt}
\begin{center}
\begin{tabular}{@{}lcccc@{}}
\toprule
\textbf{Model} & \textbf{Size} & \textbf{Dimensions} & \textbf{Quality} & \textbf{Cost} \\
\midrule
OpenAI text-embedding-3-small & API-only & 1536 & Excellent & Per API call \\
OpenAI text-embedding-ada-002 & API-only & 1536 & Good      & Per API call \\
all-MiniLM-L6-v2              & 90\,MB   & 384  & Good      & \textbf{Free, local} \\
all-mpnet-base-v2             & 420\,MB  & 768  & Better    & Free, local \\
\bottomrule
\end{tabular}
\end{center}
\vspace{6pt}

MiniLM-L6-v2 was chosen because:
\begin{itemize}
    \item \textbf{Zero API cost} --- runs entirely on CPU, no external API calls for embeddings
    \item \textbf{Small size} --- only 90\,MB, downloads once and cached locally
    \item \textbf{Fast} --- can embed 3000+ functions in under 30 seconds on a MacBook
    \item \textbf{Good quality} --- 384-dimensional vectors, sufficient for code retrieval
    \item \textbf{Student project constraint} --- OpenAI embeddings cost money; this is free
\end{itemize}

\subsubsection{How Embedding Works in Code}
The \texttt{embedder.py} file loads the model once at startup:

\begin{lstlisting}[language=Python]
from sentence_transformers import SentenceTransformer

# loads 90MB model, cached after first download
model = SentenceTransformer('all-MiniLM-L6-v2')

def embed_text(text: str) -> list[float]:
    return model.encode(text).tolist()  # returns list of 384 floats

def embed_batch(texts: list[str]) -> list[list[float]]:
    # batch processing - much faster than one at a time
    return model.encode(texts, batch_size=64, show_progress_bar=True).tolist()
\end{lstlisting}

The model loads once at service startup. Every query and every chunk uses the same model. \textbf{This is critical} --- you must use the same model for both indexing and querying, otherwise the vectors are not comparable.

\subsection{Vector Database: ChromaDB}

\subsubsection{What is a Vector Database?}
A vector database stores embedding vectors alongside their original text and metadata. Its primary operation is \textbf{similarity search}: given a query vector, find the stored vectors that are most mathematically similar (using cosine similarity or L2 distance).

\subsubsection{Why ChromaDB?}
ChromaDB was chosen over alternatives like Pinecone, Weaviate, or FAISS for these reasons:
\begin{itemize}
    \item \textbf{Local persistent storage} --- data saved to disk in \texttt{./chroma\_store/}, survives service restarts
    \item \textbf{No cloud account needed} --- Pinecone requires signup and has limits on free tier
    \item \textbf{Simple Python API} --- minimal boilerplate
    \item \textbf{Collection-based} --- each repository gets its own isolated collection
\end{itemize}

\subsubsection{Caching Strategy}
Once a repository is indexed, its ChromaDB collection persists on disk. The next time the same repository URL is analyzed, the system detects the existing collection and skips re-indexing. This means:
\begin{itemize}
    \item First analysis of \texttt{expressjs/express}: 45--90 seconds (clone + parse + embed)
    \item All subsequent analyses: under 2 seconds (use cached collection)
\end{itemize}

\subsection{Chunking Strategy: What We Index and How}

This is one of the most important design decisions in a RAG system. How you break content into chunks directly determines retrieval quality.

\subsubsection{Naive Chunking (What We Avoided)}
The simplest approach is to split files by character count (e.g., every 500 characters is one chunk). This is bad for code because:
\begin{itemize}
    \item A function might be split in half between two chunks
    \item No semantic meaning is preserved
    \item The chunk has no context about which file it came from
\end{itemize}

\subsubsection{AST-Based Function-Level Chunking (What We Use)}
We parse every JavaScript/TypeScript file using Babel's Abstract Syntax Tree (AST) parser and extract complete functions as individual chunks. Each function is one chunk --- never split.

The AST is a structured tree representation of code. For example:
\begin{lstlisting}[language=JavaScript]
function greet(name) {
    return "hello " + name;
}
\end{lstlisting}

\noindent becomes the AST node:
\begin{lstlisting}[language=Python]
{
  "type": "FunctionDeclaration",
  "id": { "name": "greet" },
  "params": [{"name": "name"}],
  "body": { ... }
}
\end{lstlisting}

We use \texttt{@babel/parser} and \texttt{@babel/traverse} in Node.js to walk this tree and extract four types of functions:
\begin{enumerate}
    \item \textbf{FunctionDeclaration} --- \texttt{function foo() \{\}}
    \item \textbf{FunctionExpression} --- \texttt{const foo = function() \{\}}
    \item \textbf{ArrowFunctionExpression} --- \texttt{const foo = () => \{\}}
    \item \textbf{ClassMethod} --- methods inside ES6 classes
\end{enumerate}

\subsubsection{Upgrade 1: Enriched Chunks}
A raw function chunk has poor retrieval quality because it lacks context. We upgraded to \textbf{enriched chunks} that include:

\begin{lstlisting}
File: lib/router/index.js
Folder: lib/router
File imports: finalhandler, debug, Layer, Route
Other functions in this file: proto.handle, proto.use, proto.route
Function: next
Type: FunctionDeclaration
Lines: 136-185

Code:
function next(err) {
    // ... actual function code
}
\end{lstlisting}

Adding the file path, folder, imports, and sibling function names dramatically improves semantic matching because the embedding model now has rich context about where this function lives in the system.

\subsubsection{Upgrade 2: File Summary Chunks}
In addition to one chunk per function, we add \textbf{one summary chunk per file}. This chunk describes the file holistically:

\begin{lstlisting}
FILE SUMMARY
File: lib/router/index.js
File name: index.js
Folder: lib/router
Imports: finalhandler, debug, Layer, Route
Total functions: 12
Named functions: proto.handle, proto.use, proto.route, proto.process_params, next
This file is part of the lib/router module.
\end{lstlisting}

This solves a critical retrieval gap. When a user asks an architectural question like \textit{``how does middleware flow work?''}, they get a file summary chunk that says \textit{``this file handles middleware dispatch''} rather than a raw function chunk. File summaries are indexed alongside function chunks, giving retrieval more surface area to match against.

\subsection{Hybrid Retrieval: BM25 + Vector Search + RRF}

\subsubsection{Why Pure Vector Search is Not Enough}
Vector search excels at semantic similarity. If you ask \textit{``how does the authentication flow work?''} it finds chunks about JWT tokens, session management, password hashing --- even if those exact words don't appear in the code.

But vector search has a critical weakness: \textbf{exact keyword matching}. If a user asks \textit{``where is authenticateToken called?''} --- this is a keyword question, not a semantic one. Vector search might miss the exact function because there are no semantically similar phrases in the chunk text.

\subsubsection{BM25: Keyword Search}
BM25 (Best Match 25) is a classical information retrieval algorithm used by search engines like Elasticsearch. It scores documents based on:
\begin{itemize}
    \item \textbf{Term frequency} --- how often the query word appears in the document
    \item \textbf{Inverse document frequency} --- rare words across the corpus get higher weight
    \item \textbf{Document length normalization} --- longer documents don't automatically score higher
\end{itemize}

BM25 excels at exact keyword matching. If you search for \texttt{authenticateToken}, BM25 will find every chunk containing that exact string, regardless of semantic similarity.

\subsubsection{Reciprocal Rank Fusion (RRF)}
Running both searches gives two ranked lists of results. We merge them using RRF:

\begin{equation}
    \text{RRF\_score}(d) =
    \frac{1}{k + \text{rank}_{\text{vector}}(d)} +
    \frac{1}{k + \text{rank}_{\text{BM25}}(d)}
\end{equation}

Where $k = 60$ is a constant from the original RRF paper. A document appearing at rank 1 in both lists gets the highest combined score. A document only in one list still gets partial credit. This is better than simply averaging scores because it doesn't require score normalization across two different scoring systems.

\subsubsection{Why Hybrid Search Outperforms Either Alone}
\begin{itemize}
    \item \textbf{Vector only:} Misses exact function names, variable names, error codes
    \item \textbf{BM25 only:} Misses semantic synonyms (``authentication'' vs.\ ``login'' vs.\ ``sign-in'')
    \item \textbf{Hybrid:} Gets the best of both --- semantic understanding \textit{and} keyword precision
\end{itemize}

Research consistently shows hybrid retrieval improves recall by 20--30\% over pure vector search for technical queries.

\subsection{The LLM: Google Gemini 2.5 Flash}

\subsubsection{Why Gemini 2.5 Flash?}
We use \texttt{gemini-2.5-flash} via Google's \texttt{google-generativeai} SDK. The reasons for this choice:

\begin{itemize}
    \item \textbf{Free tier} --- Google AI Studio provides a generous free tier for \texttt{gemini-2.5-flash}, essential for a student project
    \item \textbf{Speed} --- Flash models are optimized for low latency. Answers arrive in 3--7 seconds
    \item \textbf{Code understanding} --- Gemini models perform well on code-related tasks
    \item \textbf{Large context window} --- can handle the 5--8 retrieved chunks comfortably
    \item \textbf{Reasoning capability} --- Gemini 2.5 has strong multi-step reasoning, good for tracing execution flows across multiple files
\end{itemize}

\subsubsection{Why Not GPT-4?}
GPT-4 is excellent but:
\begin{itemize}
    \item No meaningful free tier --- requires paid OpenAI account
    \item Higher latency on complex prompts
    \item For this use case (code Q\&A with retrieved context), Gemini 2.5 Flash performs comparably
\end{itemize}

\subsubsection{The Prompt Engineering \textbf{[UPDATED]}}
The prompt evolved significantly for v1.0. Earlier versions cast the model as a ``code navigator'' that cited line numbers in prose. The current prompt makes two larger demands: it adopts an \textbf{architecture-first} teaching voice, and it asks the model to return a \textbf{single structured JSON object} containing both the prose answer and a machine-readable execution-flow graph.

\begin{lstlisting}[language=Python]
prompt = f"""You are a Senior Software Architect and Guide for this repository.

Return a SINGLE valid JSON object with EXACTLY these keys:
- "answer": a Markdown string -- the conceptual explanation.
- "execution_flow": an object with "nodes" and "edges" describing
  the runtime execution path.

answer rules:
- Explain the WHY (architecture, purpose, flow). The user sees the
  WHAT in the File Viewer and the WHERE in the Execution Flow graph.
- Architecture-first: lead with the concept, not the implementation.
- Use inline backticks for all files/functions/APIs.
- Fenced code only for short snippets (<= 20 lines); never whole files.
- NEVER mention "chunks", "embeddings", "vector search", or "retrieval".

execution_flow rules:
- Model the execution PATH that answers the question, like a whiteboard
  sketch -- not a list of retrieved files.
- Each node: id, label, type (one of: user_action, api_endpoint,
  controller, function, background_process, ai_service, database,
  output), and -- when backed by code -- file, functionName, lines.
- Edges express order; branches may split and merge.

Context Code:
{context_text}

User Question: {question}
"""
\end{lstlisting}

Key design decisions in this prompt:
\begin{itemize}
    \item \textbf{Single structured response} --- one Gemini call returns both the explanation and the graph, so the frontend never has to parse a graph out of markdown
    \item \textbf{Architecture-first voice} --- the answer explains \textit{why} the system is built this way; the implementation lives in the File Viewer
    \item \textbf{Separation of concerns} --- the prompt tells the model the user can see the ``what'' (File Viewer) and ``where'' (Execution Flow) elsewhere, so the answer stays conceptual
    \item \textbf{Hides RAG internals} --- forbidding mentions of chunks, embeddings, and retrieval makes answers read like a senior engineer, not a search engine
    \item \textbf{Typed execution graph} --- the eight node types let the frontend render each step with a distinct visual style
\end{itemize}

\begin{mdframed}[style=notebox]
\textbf{Why JSON mode?} The response is requested with \texttt{response\_mime\_type = "application/json"}, so Gemini returns parseable JSON directly. A sanitizer drops malformed nodes/edges, and if the model returns no usable graph, a deterministic fallback flow is synthesized from the retrieved functions --- the Execution Flow is therefore never empty.
\end{mdframed}

\subsection{The Complete Indexing Pipeline}

Step by step, what happens when a repository is first analyzed:

\begin{enumerate}
    \item Node.js receives \texttt{POST /api/repo/analyze} with \texttt{repoUrl}
    \item \texttt{simple-git} clones the repository into \texttt{temp\_repos/repoName/}
    \item \texttt{fileTree.js} recursively walks the directory, building a JSON tree (ignoring \texttt{node\_modules}, \texttt{.git}, \texttt{dist}, etc.)
    \item \texttt{extractFunctions.js} reads every \texttt{.js/.jsx/.ts/.tsx} file, parses it with Babel AST, extracts all functions with name, code, startLine, endLine, type
    \item \texttt{gitEvolution.js} reads the last 20 git commits, parses diffs, and builds an \texttt{evolutionMap} keyed by \texttt{filePath::functionName} \textbf{[NEW]}
    \item The extracted functions array is sent via HTTP to \texttt{POST /index} on the Python service
    \item Python service's \texttt{vector\_store.py} builds enriched chunk text for each function (file path + imports + neighbor names + code)
    \item A file summary chunk is prepended for each file
    \item All chunk texts are passed to \texttt{embed\_batch()} --- MiniLM encodes them in batches of 64
    \item All chunks + embeddings + metadata are stored in ChromaDB under a collection named after the repository
    \item Node.js returns \texttt{repoName}, \texttt{fileCount}, \texttt{totalFunctions}, and \texttt{tree} to the frontend
\end{enumerate}

\subsection{The Complete Query Pipeline}

Step by step, what happens when a user asks a question:

\begin{enumerate}
    \item Frontend sends \texttt{POST /api/repo/query} with \texttt{repoName} and \texttt{question}
    \item Node.js proxies this to \texttt{POST /query} on the Python service
    \item Python's \texttt{query\_collection()} runs \textbf{two searches simultaneously}:
    \begin{itemize}
        \item \textbf{Vector search:} Embeds the question with MiniLM, queries ChromaDB for top 20 semantically similar chunks
        \item \textbf{BM25 search:} Fetches all documents from ChromaDB, builds in-memory BM25 index, scores all chunks against the question tokens, takes top 20
    \end{itemize}
    \item RRF merges the two ranked lists into one combined ranking
    \item Chunks from test files (\texttt{test/}, \texttt{tests/}, \texttt{examples/}, \texttt{.spec.}) are deprioritized
    \item Top 5 chunks passed to \texttt{ask\_llm()} in \texttt{llm.py}
    \item Gemini 2.5 Flash receives the structured prompt and returns a JSON object: a markdown \texttt{answer} plus an \texttt{execution\_flow} graph \textbf{[UPDATED]}
    \item Execution-flow node line ranges are reconciled against the indexed function metadata (\texttt{get\_function\_index}) so every code node points at real lines \textbf{[NEW]}
    \item Source files and source functions extracted from chunk metadata
    \item Full response (answer + sources + execution flow) returned to Node.js, forwarded to frontend \textbf{[UPDATED]}
\end{enumerate}

\newpage

% ═══════════════════════════════════════════════════════
%  SECTION 4: NODE.JS BACKEND
% ═══════════════════════════════════════════════════════
\section{Node.js Backend: Detailed Breakdown}

The backend lives in the \texttt{backend/} folder and runs on Express.js port 8000. The backend has been refactored into an MVC structure: routing logic lives in \texttt{backend/routes/repo.js} and all business logic has been moved to \texttt{backend/controllers/repoController.js}. \textbf{[NEW]}

\subsection{Key Libraries and Their Roles}

\subsubsection{Express.js}
The web framework. Handles HTTP routing, middleware, request/response lifecycle. All API endpoints are defined here.

\subsubsection{simple-git}
A Node.js wrapper around Git. Used to clone repositories programmatically:
\begin{lstlisting}[language=JavaScript]
await simpleGit().clone(repoUrl, clonePath)
\end{lstlisting}
This is equivalent to running \texttt{git clone} from the command line but within the Node.js process. Cloned repos are stored in \texttt{temp\_repos/} with the repo name as the folder name. \texttt{simple-git} is also used inside \texttt{gitEvolution.js} to run \texttt{git show} commands for commit diff parsing. \textbf{[NEW]}

\subsubsection{@babel/parser and @babel/traverse}
These are the AST parsing libraries. \texttt{@babel/parser} reads JavaScript/TypeScript source code and produces an Abstract Syntax Tree. \texttt{@babel/traverse} walks that tree and calls your callback for each node type.

The parser is configured to handle all modern JS/TS syntax:
\begin{lstlisting}[language=JavaScript]
ast = parser.parse(code, {
    sourceType: 'unambiguous',  // handles both ESM and CJS
    plugins: [
        'jsx',              // React JSX syntax
        'typescript',       // TypeScript types
        'classProperties',  // class fields
        'optionalChaining', // ?. operator
        'nullishCoalescingOperator' // ?? operator
    ],
    errorRecovery: true  // don't crash on syntax errors
})
\end{lstlisting}

\subsubsection{axios}
HTTP client used by Node.js to call the Python AI service. All communication between Node.js and Python goes through HTTP:
\begin{lstlisting}[language=JavaScript]
const response = await axios.post('http://localhost:8001/index', {
    repo_name: repoName,
    extracted_functions: extractedFunctions
}, { timeout: 300000 })
\end{lstlisting}

\subsubsection{cors}
Middleware that allows the React frontend (running on port 5173) to make requests to the Express backend (running on port 8000). Without CORS headers, browsers block cross-origin requests.

\subsubsection{dotenv}
Loads environment variables from the \texttt{.env} file so configuration such as the server \texttt{PORT} and the CORS allow-list is not hardcoded in source code. \textbf{[UPDATED]}

\subsection{API Endpoints}

\subsubsection{POST /api/repo/analyze}
\textbf{Purpose:} Clone a repository, extract functions, compute git evolution, trigger indexing\\[2pt]
\textbf{Body:} \texttt{\{ repoUrl: "https://github.com/..." \}}\\[4pt]
\textbf{Flow:}
\begin{enumerate}
    \item Validate URL starts with \texttt{https://github.com}
    \item Derive repo name from URL (\texttt{owner/repo} $\rightarrow$ \texttt{owner\_repo})
    \item Check if already cloned in \texttt{temp\_repos/} (cache check)
    \item If not cached: \texttt{simpleGit().clone(url, path)}
    \item Build file tree (ignoring \texttt{node\_modules}, \texttt{.git}, etc.)
    \item Count files, reject if over 300
    \item Extract functions using Babel AST
    \item Compute git evolution map via \texttt{gitEvolution.js} \textbf{[NEW]}
    \item POST extracted functions to Python \texttt{/index}
    \item Return \texttt{\{repoName, fileCount, totalFunctions, tree\}} to frontend
\end{enumerate}

\subsubsection{POST /api/repo/query}
\textbf{Purpose:} Proxy a natural language question to the Python AI service\\[2pt]
\textbf{Body:} \texttt{\{ repoName: "...", question: "..." \}}\\[4pt]
\textbf{Flow:}
\begin{enumerate}
    \item Validate repoName and question present
    \item POST to Python \texttt{/query} with \texttt{repo\_name} and \texttt{question}
    \item Return Python's response directly to frontend
\end{enumerate}

\subsubsection{GET /api/repo/file}
\textbf{Purpose:} Read and return the content of a specific file from the cloned repo\\[2pt]
\textbf{Query params:} \texttt{repoName}, \texttt{filePath}\\[4pt]
\textbf{Flow:}
\begin{enumerate}
    \item Construct full path: \texttt{temp\_repos/repoName/filePath}
    \item Check file exists
    \item Read file with \texttt{fs.readFileSync}
    \item Return \texttt{\{filePath, content, lines, totalLines\}}
\end{enumerate}

This endpoint powers the File Viewer, which can be opened by clicking a file in the tree, a citation in an answer, or a node in the Execution Flow. \textbf{[UPDATED]}

\subsubsection{POST /api/repo/evolution \textbf{[NEW]}}
\textbf{Purpose:} Return the git change history for functions in the repository, with AI-generated summaries\\[2pt]
\textbf{Body:} \texttt{\{ repoName: "..." \}}\\[4pt]
\textbf{Flow:}
\begin{enumerate}
    \item Look up the cloned repo in \texttt{temp\_repos/repoName/}
    \item Call \texttt{getCodeEvolution(repoPath, extractedFunctions)} from \texttt{gitEvolution.js}
    \item Receive \texttt{evolutionMap} keyed as \texttt{filePath::functionName}
    \item POST evolution data to Python \texttt{/evolution-summary}
    \item Return per-function change history and summaries to frontend
\end{enumerate}

\subsection{The File Tree Builder (fileTree.js)}

Recursively walks a directory and returns a JSON tree:
\begin{lstlisting}[language=JavaScript]
const IGNORE = new Set([
    'node_modules', '.git', '.next', 'dist',
    'build', '.cache', 'coverage'
])

function buildFileTree(dirPath, relativePath = '') {
    const entries = fs.readdirSync(dirPath, { withFileTypes: true })
    const result = []
    
    for (const entry of entries) {
        if (IGNORE.has(entry.name)) continue
        // recursively build children for folders
        // add file nodes for files
    }
    return result
}
\end{lstlisting}

The IGNORE set is critical --- without it, \texttt{node\_modules} alone would produce thousands of nodes and crash the indexing.

\subsection{The Function Extractor (extractFunctions.js)}

For each \texttt{.js/.jsx/.ts/.tsx} file, this utility:
\begin{enumerate}
    \item Reads the file content
    \item Skips files over 500\,KB (likely generated/minified code)
    \item Parses with Babel to get AST
    \item Traverses AST looking for 4 node types
    \item For each function found, stores: \texttt{name}, \texttt{type}, \texttt{code} (sliced from source), \texttt{startLine}, \texttt{endLine}
    \item Returns array of \texttt{\{filePath, functionCount, functions\}}
\end{enumerate}

\subsection{The Git Evolution Tracker (gitEvolution.js) \textbf{[NEW]}}

This new utility computes a function-level change history by combining git diff output with the already-extracted function line ranges.

\textbf{How it works:}
\begin{enumerate}
    \item Reads the last 20 commits from the cloned repo using \texttt{simple-git}
    \item For each commit, runs \texttt{git show --unified=0} to get a minimal diff (no surrounding context lines)
    \item Parses the diff to extract: which file was changed, which line numbers were added (\texttt{+}) and deleted (\texttt{-}), and the commit message + timestamp
    \item For each changed line, checks which extracted function's \texttt{startLine}--\texttt{endLine} range it falls within
    \item Maps the change to that function, building an \texttt{evolutionMap}
\end{enumerate}

\textbf{The evolutionMap structure:}
\begin{lstlisting}[language=JavaScript]
{
  "lib/router/index.js::proto.handle": [
    {
      commitHash: "a3f9b12",
      message:   "fix: handle async errors in middleware",
      timestamp: "2024-11-03T14:22:00Z",
      linesAdded: 3,
      linesDeleted: 1
    },
    // ... more commits
  ],
  "middleware/auth.js::verifyToken": [ ... ]
}
\end{lstlisting}

This map is then sent to the Python AI service, which asks Gemini to produce a human-readable summary of how each function has evolved over the commits. The summary describes things like ``this function was refactored to handle async errors in November and had its JWT logic simplified in December.''

\newpage

% ═══════════════════════════════════════════════════════
%  SECTION 5: PYTHON AI SERVICE
% ═══════════════════════════════════════════════════════
\section{Python AI Service: File-by-File Breakdown}

The service lives in the \texttt{service/} folder and runs FastAPI on port 8001.

\subsection{embedder.py}
Single responsibility: load and expose the MiniLM model. The model loads once at module import time (when FastAPI starts). All other modules import from here.

\textbf{Key functions:}
\begin{itemize}
    \item \texttt{embed\_text(text)} --- embed a single string (used for queries)
    \item \texttt{embed\_batch(texts)} --- embed a list of strings efficiently (used during indexing)
\end{itemize}

\subsection{vector\_store.py}
All ChromaDB operations are isolated here. Nothing else touches the database directly.

\textbf{Key functions:}
\begin{itemize}
    \item \texttt{get\_or\_create\_collection(repo\_name)} --- get existing or create new ChromaDB collection
    \item \texttt{is\_already\_indexed(repo\_name)} --- check if collection exists and has documents (caching check)
    \item \texttt{delete\_collection(repo\_name)} --- force re-index by deleting existing collection
    \item \texttt{\_extract\_file\_metadata(file\_data)} --- extract import statements and all function names from a file (used for enriching chunks)
    \item \texttt{\_build\_chunk\_text(file\_path, func, file\_metadata)} --- build the enriched text for one function chunk
    \item \texttt{\_build\_file\_summary\_chunk(file\_path, file\_metadata, file\_data)} --- build the summary text for one file
    \item \texttt{index\_functions(repo\_name, extracted\_functions)} --- main indexing function: builds all chunks, embeds them in batch, stores in ChromaDB
    \item \texttt{query\_collection(repo\_name, question, top\_k)} --- hybrid search: BM25 + vector + RRF, returns top chunks
    \item \texttt{\_tokenize(text)} --- simple regex tokenizer for BM25 (extracts word tokens from code text)
    \item \texttt{\_reciprocal\_rank\_fusion(vector\_ranking, bm25\_ranking)} --- merges two ranked lists using RRF formula
\end{itemize}

\subsection{llm.py}
Single responsibility: call Gemini and return a structured response.

\textbf{Key function:} \texttt{ask\_llm(question, context\_chunks)}\\[3pt]
Takes the question and retrieved chunks, builds the structured prompt, calls Gemini, extracts \texttt{source\_files} and \texttt{source\_functions} from chunk metadata, returns everything.

The \texttt{client} object is also exported and reused by \texttt{main.py} for the evolution summary endpoint, so Gemini is only initialized once. \textbf{[NEW]}

\subsection{main.py}
FastAPI entry point. Routes only --- no business logic here.

\textbf{Endpoints:}
\begin{itemize}
    \item \texttt{GET /health} --- health check
    \item \texttt{POST /index} --- receive extracted functions from Node.js, index them
    \item \texttt{POST /query} --- receive question + repo name, return answer + sources
    \item \texttt{DELETE /index/\{repo\_name\}} --- force delete collection for re-indexing
    \item \texttt{POST /evolution-summary} --- receive evolutionMap from Node.js, ask Gemini to summarize per-function change history, return summaries \textbf{[NEW]}
\end{itemize}

\newpage

% ═══════════════════════════════════════════════════════
%  SECTION 6: FRONTEND
% ═══════════════════════════════════════════════════════
\section{Frontend: React Application}

The frontend lives in \texttt{frontend/} and runs on Vite + React, port 5173.

\subsection{Technology Stack}
\begin{itemize}
    \item \textbf{React 18} --- component-based UI
    \item \textbf{Vite} --- build tool and dev server (faster than Create React App)
    \item \textbf{React Router v6} --- client-side navigation between pages
    \item \textbf{Tailwind CSS v3} --- utility-first styling
    \item \textbf{axios} --- HTTP client for API calls
    \item \textbf{react-markdown} --- renders LLM responses as formatted markdown
    \item \textbf{React Flow (@xyflow/react) + Dagre} --- renders the Execution Flow graph with automatic layout \textbf{[NEW]}
    \item \textbf{Shiki} --- syntax highlighting for the File Viewer and code blocks \textbf{[NEW]}
\end{itemize}

\subsection{Pages}

\subsubsection{LandingPage.jsx}
The entry point. Contains:
\begin{itemize}
    \item Hero section with headline and subheadline
    \item GitHub URL input bar with Analyze button
    \item Step-by-step loading progress panel (appears while analyzing)
    \item Three feature cards (Zero Setup, AI Q\&A, Focus Map)
    \item How it works section with 3 numbered steps
\end{itemize}

On clicking Analyze:
\begin{enumerate}
    \item Calls \texttt{POST /api/repo/analyze} with 5-minute timeout (cloning + indexing takes time)
    \item Shows loading panel with pipeline steps
    \item On success: navigates to \texttt{/chat} passing \texttt{repoName}, \texttt{fileCount}, \texttt{totalFunctions}, \texttt{tree} via React Router state
    \item On error: shows error message inline
\end{enumerate}

\subsubsection{ChatPage.jsx}
The main application view. Split layout --- left 60\%, right 40\%.

\textbf{Left panel:}
\begin{itemize}
    \item Top bar with repo name, file count badge, functions count badge, back button
    \item Scrollable message history
    \item Bouncing dots loading indicator while waiting for AI
    \item Fixed input bar at bottom (textarea + send button)
    \item Enter key submits, Shift+Enter for newline
\end{itemize}

\textbf{Right panel (3 tabs):}
\begin{itemize}
    \item \textbf{File Tree} --- collapsible folder/file navigator
    \item \textbf{Execution Flow} --- interactive graph of the execution path \textbf{[UPDATED]}
    \item \textbf{File Viewer} --- code viewer with line numbers
\end{itemize}

\textbf{State management:}
\begin{lstlisting}[language=JavaScript]
const [messages, setMessages]   = useState(...)  // chat history (persisted)
const [activeTab, setActiveTab] = useState(...)  // right panel tab
const [focusData, setFocusData] = useState(...)  // execution flow shown
const [selectedFile, setSelectedFile]   = useState(...) // file viewer
const [highlight, setHighlight]         = useState(...) // lines to highlight
const [selectedFunction, setSelectedFunction] = useState(...) // for Git Evolution
const [evolutionTarget, setEvolutionTarget]   = useState(...) // drawer target
\end{lstlisting}

\textbf{State persistence \textbf{[NEW]}:} On every meaningful change, the chat history, selected repository and metadata, per-answer execution flows, citations, and File Viewer state are written (debounced) to \texttt{localStorage} under a single versioned key via \texttt{lib/chatStorage.js}. On load, \texttt{ChatPage} hydrates from router state for a fresh analysis, or from LocalStorage on a refresh --- so the conversation is restored automatically with no ``Restore'' button. Corrupted or out-of-version data is cleared gracefully rather than crashing the app.

\subsection{Components}

\subsubsection{ChatMessage.jsx}
Renders a single chat bubble. User messages are right-aligned with indigo background. AI messages are left-aligned with dark background.

AI messages use \texttt{react-markdown} with custom component overrides to render:
\begin{itemize}
    \item \texttt{**bold**} as white bold text
    \item Bullet lists with proper spacing
    \item Numbered lists for sequential steps
    \item Inline code (\texttt{`code`}) with purple background highlight
\end{itemize}

Each AI message whose response includes an execution flow shows a \texttt{"View Execution Flow →"} button that renders the graph in the right panel. \textbf{[UPDATED]} In addition, file and function references inside the answer (rendered as inline code) are matched against the message's \texttt{sourceFiles} and \texttt{sourceFunctions} and become \textbf{clickable citations}: clicking a file opens it in the File Viewer; clicking a function opens its file, scrolls to it, and briefly highlights its lines. Unmatched code spans stay inert, avoiding false-positive links. \textbf{[NEW]}

\subsubsection{FileTree.jsx}
Recursive tree component. Each node is either a folder (clickable to expand/collapse) or a file (clickable to open in File Viewer).

\begin{lstlisting}[language=JavaScript]
function TreeNode({ node, onFileClick }) {
    const [open, setOpen] = useState(true)
    
    if (node.type === 'file') {
        return <div onClick={() => onFileClick(node.path)}>...</div>
    }
    // folder with collapsible children
}
\end{lstlisting}

Clicking a file calls \texttt{onFileClick(filePath)}, which sets \texttt{selectedFile} in ChatPage and switches to the File Viewer tab.

\subsubsection{FileViewer.jsx \textbf{[UPDATED]}}
Fetches and displays file content via the \texttt{GET /api/repo/file} endpoint, rendered through a shared \texttt{CodeBlock} component that applies \textbf{Shiki syntax highlighting} with line numbers.

When opened from a citation or an Execution Flow node, the viewer scrolls to the referenced line range and briefly highlights it. If the open file corresponds to a selected function, a contextual \texttt{"View Git Evolution"} action appears in the header.

\subsubsection{FocusMap.jsx \textbf{[UPDATED]}}
Originally a grouped list of retrieved files and functions, this component was redesigned into an \textbf{Execution Flow} graph. It no longer reconstructs anything from the answer text --- it renders the structured \texttt{execution\_flow} object (\texttt{nodes} + \texttt{edges}) returned by the AI service.

\begin{itemize}
    \item \textbf{React Flow + Dagre} --- nodes and edges are laid out automatically with Dagre (top-to-bottom), producing a workflow diagram that branches and merges rather than a vertical list
    \item \textbf{Typed nodes} --- each node is styled by its type (user action, API endpoint, controller, function, background process, AI service, database, output), so the execution path is easy to read at a glance
    \item \textbf{Interactive navigation} --- clicking a code-backed node opens its file in the File Viewer, scrolls to the function, and highlights it; conceptual nodes (e.g.\ ``Clone Repository'') are non-navigable
    \item \textbf{Resilience} --- malformed nodes/edges are filtered, and an empty graph falls back to the flow synthesized by the AI service
\end{itemize}

The graph answers ``what happens after this?'' for the user's specific question, rather than ``which files were retrieved?''

\subsubsection{GitEvolution.jsx \textbf{[NEW]}}
A right-side drawer that surfaces the git evolution feature on the frontend. When a function is selected (from an Execution Flow node), a \texttt{"View Git Evolution"} action opens this drawer, which calls \texttt{GET /api/repo/evolution} and shows a commit timeline for that function together with the Gemini-generated change summary. If the function has no recorded history, it shows a graceful empty state rather than failing. (The earlier \texttt{RepoHeader.jsx} placeholder was removed; repository metadata badges are rendered inline in the \texttt{ChatPage} top bar.)

\newpage

% ═══════════════════════════════════════════════════════
%  SECTION 7: PHASE-WISE DEVELOPMENT
% ═══════════════════════════════════════════════════════
\section{Phase-Wise Development: What Was Built, Issues Found, How Fixed}

\subsection{Phase 1: Core Pipeline --- Clone, Parse, Index}

\subsubsection{What was built}
\begin{enumerate}
    \item \textbf{POST /api/repo/analyze endpoint} --- accepts GitHub URL, clones with simple-git, builds file tree
    \item \textbf{fileTree.js} --- recursive directory walker with ignore list
    \item \textbf{extractFunctions.js} --- Babel AST parser for JS/TS function extraction
    \item \textbf{Python /index endpoint} --- receives functions, embeds with MiniLM, stores in ChromaDB
    \item \textbf{Python /query endpoint} --- basic vector-only search, sends to Gemini, returns answer
    \item \textbf{POST /api/repo/query in Node.js} --- proxy to Python
\end{enumerate}

\subsubsection{Issues Found}

\textbf{Issue 1: Port 5000 blocked on Mac}\\
macOS Monterey and later reserves port 5000 for AirPlay Receiver. Every request to port 5000 got a 403 Forbidden --- never reached Express. Fixed by switching to port 8000.

\vspace{4pt}
\textbf{Issue 2: Module system mismatch}\\
Initial code mixed ES modules (\texttt{import/export}) with CommonJS (\texttt{require/module.exports}). Node.js cannot mix these. Fixed by standardizing all backend files to CommonJS.

\vspace{4pt}
\textbf{Issue 3: ChromaDB include parameter}\\
ChromaDB's \texttt{query()} method does not accept \texttt{"ids"} in the include list (IDs are always returned automatically). Caused \texttt{ValueError} on every query. Fixed by removing \texttt{"ids"} from the include list.

\vspace{4pt}
\textbf{Issue 4: Gemini SDK deprecation}\\
The \texttt{google-generativeai} Python package was deprecated by Google. Updated to \texttt{google-genai} and changed the API call pattern to:
\begin{lstlisting}[language=Python]
client = genai.Client(api_key=os.getenv("GEMINI_API_KEY"))
response = client.models.generate_content(
    model="gemini-2.5-flash",
    contents=prompt
)
\end{lstlisting}
Note: During environment setup (Phase 6), the virtual environment was found to have \texttt{google-generativeai} installed rather than \texttt{google-genai}. The code was temporarily reverted to \texttt{google-generativeai} to match the installed package. Migration back to \texttt{google-genai} is pending. A \texttt{FutureWarning} is emitted at service startup but does not affect functionality. \textbf{[UPDATED]}

\vspace{4pt}
\textbf{Issue 5: numpy binary incompatibility on Mac}\\
After upgrading numpy, scikit-learn became incompatible (\texttt{numpy.dtype size changed}). Fixed by: \texttt{pip uninstall numpy scikit-learn -y} and \texttt{pip install numpy scikit-learn}.

\vspace{4pt}
\textbf{Issue 6: Sending 22,000 lines of JSON via Postman}\\
Attempted to manually POST the extracted functions to the Python service from Postman --- impossible due to size. Fixed by wiring Node.js to automatically call Python's \texttt{/index} after extraction (server-to-server call with 5 minute timeout).

\subsubsection{Result at End of Phase 1}
Full pipeline working: paste GitHub URL $\rightarrow$ functions extracted $\rightarrow$ indexed in ChromaDB $\rightarrow$ questions answered via Postman. 3093 functions from \texttt{expressjs/express} indexed in \textasciitilde30 seconds.

\subsection{Phase 2: RAG Quality Improvements}

\subsubsection{Problem Identified}
The initial RAG worked but answers were poor quality. Root causes:
\begin{enumerate}
    \item Raw function code has no context --- the embedding model doesn't know which file or module the function belongs to
    \item 70\% of functions are named ``anonymous'' --- no semantic signal in the chunk text
    \item Pure vector search misses exact keyword matches
    \item Test files dominated retrieval results --- \texttt{test/} folder had 80+ files in Express repo, overwhelming source files
\end{enumerate}

\subsubsection{Upgrade 1: Enriched Chunks}
Built \texttt{\_build\_chunk\_text()} in \texttt{vector\_store.py} that prepends file path, folder context, imports, and sibling function names to each chunk before embedding. This gave the embedding model rich semantic context. Retrieval quality improved significantly --- source files started appearing in results instead of only test files.

\subsubsection{Upgrade 2: File Summary Chunks}
Added \texttt{\_build\_file\_summary\_chunk()} that generates one holistic summary per file. For a 200-file repo, this adds 200 extra chunks that answer architectural questions. Total chunks went from 3093 to 3227 for Express.

\subsubsection{Upgrade 3: Hybrid Search (BM25 + Vector + RRF)}
Installed \texttt{rank-bm25} library. Rewrote \texttt{query\_collection()} to run both searches in parallel, then merge with RRF. This solved the exact keyword matching problem --- queries with specific function names, file names, or error codes now reliably retrieve the right chunks.

\subsubsection{Test File Filtering}
Added a filter in \texttt{main.py} to deprioritize chunks from \texttt{test/}, \texttt{tests/}, \texttt{examples/}, \texttt{.spec.}, \texttt{.test.} paths. If filtered results are empty, falls back to original results.

\subsubsection{Result at End of Phase 2}
Noticeably better retrieval. Auth questions return \texttt{controller/authController.js} instead of test files. File summary chunks appear in Focus Map for architectural questions.

\subsection{Phase 3: Prompt Engineering}

\subsubsection{Problem Identified}
The LLM was receiving good chunks but producing generic answers like ``Authentication appears to involve issuing a token.'' Not useful for a developer who needs to understand the system.

\subsubsection{Solution}
Rewrote the prompt with explicit structure:
\begin{itemize}
    \item \textbf{Role definition} --- expert code navigator
    \item \textbf{Step-by-step answering instructions} --- one-sentence answer, then details, trace flow, cite line numbers
    \item \textbf{Formatting rules} --- bold for file names, bullet points, numbered steps for flows
    \item \textbf{Strict hallucination prevention} --- only state what is clearly visible in chunks
    \item \textbf{Gap reporting} --- explicitly state what information is missing when chunks are insufficient
\end{itemize}

\subsubsection{Result}
Answer quality dramatically improved. Responses now include numbered sections, bold citations, execution flow tracing, and a ``What is missing'' section when chunks don't fully cover the question. This ``missing'' section is particularly valuable --- it's honest rather than hallucinating.

\subsection{Phase 4: Frontend Development}

\subsubsection{Landing Page}
Built with React + Tailwind CSS. Dark theme (\texttt{\#0a0a0a} background, white text, indigo accent). Contains URL input, loading state panel, feature cards, how-it-works section.

\subsubsection{Chat Page}
Split layout with three-tab right panel. Connected to real backend endpoints. Markdown rendering via \texttt{react-markdown}. Bouncing dots loading indicator.

\subsubsection{File Tree Component}
Recursive tree with expand/collapse folders. Clicking files triggers File Viewer tab.

\subsubsection{File Viewer Component}
Fetches file content from \texttt{GET /api/repo/file}. Renders as table with line numbers on left column. Preserves indentation with \texttt{whitespace-pre}.

\subsubsection{Focus Map Component}
Initially just showed a text list of source files. Upgraded to grouped node layout (file as container, functions as child nodes). Added hover tooltip with code preview using fixed positioning to avoid overflow clipping issues.

\subsubsection{Issues During Frontend Development}

\textbf{Issue 1: All component files empty}\\
\texttt{touch} command creates empty files. VS Code showed the file path but contents were empty. Every import caused ``does not provide default export'' error. Fixed by manually pasting and saving each component.

\vspace{4pt}
\textbf{Issue 2: Tailwind v4 incompatibility}\\
\texttt{npx tailwindcss init -p} fails on Tailwind v4 (no init command). Fixed by uninstalling v4 and installing \texttt{tailwindcss@3} which still has the \texttt{init} command.

\vspace{4pt}
\textbf{Issue 3: Hover tooltip clipped}\\
The code preview tooltip was positioned \texttt{absolute} inside an overflow-hidden container. Tooltip was invisible (clipped). Fixed by switching to \texttt{position: fixed} with JavaScript-calculated screen coordinates, keeping the tooltip in the viewport layer above all containers.

\vspace{4pt}
\textbf{Issue 4: Tooltip going off-screen}\\
Fixed-position tooltip appeared off the right edge of the screen. Fixed by checking \texttt{window.innerWidth - rect.right} (available space on right) and showing tooltip on the left side when insufficient space exists on the right.

\subsection{Phase 5: GitHub and Documentation}

\begin{enumerate}
    \item Created \texttt{.gitignore} to exclude \texttt{node\_modules}, \texttt{temp\_repos}, \texttt{chroma\_store}, \texttt{.env} files
    \item Initialized git repository
    \item Pushed to \texttt{github.com/Sahil0282/repo-explorer}
    \item Added comprehensive README with architecture diagram, setup instructions, tech stack explanation
    \item Added RepoExplorer as first project on resume with 6 bullet points covering the technically impressive aspects
\end{enumerate}

\subsection{Phase 6: Git Evolution Feature \textbf{[NEW]}}

\subsubsection{Motivation}
After completing the core RAG Q\&A pipeline, a second feature was added to answer a fundamentally different type of question: not ``how does this code work?'' but ``how has this code changed over time?'' This is useful for understanding why a function is structured the way it is, or for identifying which functions have been heavily modified (indicating instability or active development).

\subsubsection{What Was Built}
\begin{enumerate}
    \item \textbf{gitEvolution.js} in \texttt{backend/utils/} --- reads the last 20 commits using \texttt{simple-git}, runs \texttt{git show --unified=0} per commit, parses added/deleted line ranges, and maps them to extracted functions by line range overlap
    \item \textbf{POST /api/repo/evolution} in Node.js --- new endpoint that triggers the evolution computation and proxies results to the Python service
    \item \textbf{POST /evolution-summary} in Python \texttt{main.py} --- new FastAPI endpoint that receives the evolutionMap and uses the shared Gemini client to generate per-function change summaries
    \item \textbf{MVC refactor} --- took the opportunity to separate routing from business logic. All controller functions moved to \texttt{backend/controllers/repoController.js}; \texttt{backend/routes/repo.js} now only defines routes
\end{enumerate}

\subsubsection{Issues During Phase 6}

\textbf{Issue 1: Wrong uvicorn launch path}\\
Running \texttt{uvicorn main:app} from the repo root failed because \texttt{main.py} is inside \texttt{service/}. Fixed by \texttt{cd service} first, then \texttt{uvicorn main:app --reload --port 8001}.

\vspace{4pt}
\textbf{Issue 2: rank-bm25 not in requirements.txt}\\
The \texttt{rank-bm25} package was added to \texttt{vector\_store.py} during Phase 2 but was never added to \texttt{requirements.txt}. Service failed to start with \texttt{ModuleNotFoundError: No module named 'rank\_bm25'}. Fixed by \texttt{pip install rank-bm25}. The \texttt{requirements.txt} file must be regenerated with \texttt{pip freeze} to capture this dependency permanently.

\vspace{4pt}
\textbf{Issue 3: Gemini SDK mismatch}\\
\texttt{llm.py} was importing \texttt{from google import genai} (the new \texttt{google-genai} style), but the virtual environment had \texttt{google-generativeai} installed (the old package). Caused \texttt{ImportError} at startup. Temporarily resolved by reverting \texttt{llm.py} to use \texttt{google.generativeai}. Full migration to \texttt{google-genai} is a known pending task.

\subsubsection{Result at End of Phase 6}
All three services start cleanly. The evolution pipeline is operational: pasting a repo URL triggers both the RAG indexing pipeline and the git evolution computation in a single \texttt{/analyze} call. Per-function change summaries are generated by Gemini and available via the \texttt{/evolution} endpoint.

\subsection{Phase 7: v1.0 --- Execution Flow, Citations, and Persistence \textbf{[NEW]}}

\subsubsection{Motivation}
The earlier Focus Map proved the graph worked but only listed retrieved files --- it did not explain how the code actually runs. Phase 7 turned the right panel into an \textbf{Execution Flow} that a developer can read like a whiteboard sketch, and tightened the loop between explanation, visualization, and implementation.

\subsubsection{What Was Built}
\begin{enumerate}
    \item \textbf{Structured LLM output} --- the query prompt now returns a single JSON object containing both the markdown answer and an \texttt{execution\_flow} graph (nodes, edges, node types, and code locations). A sanitizer and a deterministic fallback guarantee a valid, non-empty graph.
    \item \textbf{Execution Flow (React Flow + Dagre)} --- \texttt{FocusMap.jsx} renders that graph with automatic layout and eight typed node styles. The flow is generated per question and changes with the question.
    \item \textbf{Line grounding} --- execution-flow node line ranges are reconciled against the indexed function metadata so every code node points at real lines.
    \item \textbf{Clickable citations} --- file/function references in answers resolve against the message's sources and navigate to the File Viewer with scroll and highlight.
    \item \textbf{Git Evolution frontend} --- a drawer (\texttt{GitEvolution.jsx}) exposes the existing evolution backend: selecting a function reveals a commit timeline and an AI summary, with a graceful empty state.
    \item \textbf{Local chat persistence} --- the full session (chat, repository, execution flows, citations, File Viewer state) is saved to LocalStorage and restored seamlessly on refresh.
    \item \textbf{Repository cleanup lifecycle} --- \texttt{temp\_repos/} is bounded: when a new repository is analyzed, stale clones beyond the most-recent few are removed, while the active repository and any in-flight analysis are never deleted.
\end{enumerate}

\subsubsection{Result at End of Phase 7}
A developer can ask a question, read a concise architecture-first explanation, see the exact execution path as an interactive graph, click any code node or citation to jump to the implementation, and open the git history for that function --- with the whole session surviving a page refresh.

\newpage

% ═══════════════════════════════════════════════════════
%  SECTION 8: KNOWN LIMITATIONS
% ═══════════════════════════════════════════════════════
\section{Current Limitations of the MVP}

\subsection{Repository Constraints}
\begin{itemize}
    \item \textbf{Public repos only} --- private repositories require GitHub OAuth token integration
    \item \textbf{JS/TS only} --- Python, Java, Go, Rust files are not parsed
    \item \textbf{Under 300 files} --- larger repos are rejected to control embedding cost and time
    \item \textbf{Under 50\,MB} --- very large files (generated/minified code) are skipped
\end{itemize}

\subsection{RAG Quality Limitations}
\begin{itemize}
    \item \textbf{Anonymous functions} --- 60--70\% of functions have no name (callbacks, arrow functions, IIFE). These are labeled \texttt{anonymous} or \texttt{fn @ L\{line\}} which provides poor semantic signal for retrieval
    \item \textbf{No call graph} --- the system doesn't know which function calls which. It can't trace execution flows across files unless the LLM infers it from the chunks
    \item \textbf{Semantic gap for framework-specific terminology} --- Express internally calls middleware ``next'', not ``middleware''. Queries using the domain term don't match the implementation term
    \item \textbf{Test file dominance} --- repositories with many test files (like Express with 80+ test files) have test chunks outnumbering source chunks, biasing retrieval
\end{itemize}

\subsection{Infrastructure Limitations}
\begin{itemize}
    \item \textbf{No deployment} --- currently localhost only; requires all three services running locally
    \item \textbf{Local-only session persistence} --- the chat, repository, execution flows, and File Viewer state survive a refresh via LocalStorage, but history is per-browser; there is no server-side or cross-device persistence \textbf{[UPDATED]}
    \item \textbf{No rate limiting} --- a malicious user could trigger unlimited indexing jobs
    \item \textbf{BM25 rebuilt on every query} --- fetches all ChromaDB documents and rebuilds BM25 index in memory for each query (slow for large repos; should be cached)
    \item \textbf{No auth} --- anyone who knows the localhost URL can use it
    \item \textbf{requirements.txt out of sync} --- \texttt{rank-bm25} was installed manually and not yet added to \texttt{requirements.txt}. A fresh \texttt{pip install -r requirements.txt} will fail until the file is regenerated with \texttt{pip freeze} \textbf{[NEW]}
    \item \textbf{Gemini SDK FutureWarning} --- \texttt{llm.py} currently uses \texttt{google-generativeai} (deprecated) instead of the newer \texttt{google-genai}. Service works but emits a warning at startup \textbf{[NEW]}
    \item \textbf{Git evolution only covers last 20 commits} --- repositories with long histories will have incomplete evolution data \textbf{[NEW]}
\end{itemize}

\newpage

% ═══════════════════════════════════════════════════════
%  SECTION 9: FUTURE IMPROVEMENTS
% ═══════════════════════════════════════════════════════
\section{Ongoing and Future Improvements}

\subsection{Improvement 1: Better Anonymous Function Naming (High Priority)}

\textbf{Problem:} Most Express route handlers look like:
\begin{lstlisting}[language=JavaScript]
app.get('/users', async (req, res) => {
    // handler code
})
\end{lstlisting}
This is an arrow function with no name --- stored as \texttt{anonymous}. But from context, we know its name should be \texttt{GET /users route handler}.

\textbf{Solution:} During AST traversal in \texttt{extractFunctions.js}, when an anonymous function is encountered, look at its parent node. If the parent is:
\begin{itemize}
    \item \texttt{app.get('/path', fn)} $\rightarrow$ name it \texttt{GET /path handler}
    \item \texttt{router.post('/path', fn)} $\rightarrow$ name it \texttt{POST /path handler}
    \item \texttt{app.use(middleware)} $\rightarrow$ name it \texttt{middleware: \{fn body preview\}}
    \item \texttt{mongoose.Schema(\{...\})} $\rightarrow$ name it \texttt{Schema definition}
    \item \texttt{Promise.then(fn)} $\rightarrow$ name it \texttt{then callback}
\end{itemize}

This dramatically improves both the focus map labels and the retrieval quality because the chunk text now says ``GET /users route handler'' instead of ``anonymous''.

\textbf{Implementation:} In \texttt{ArrowFunctionExpression} handler in \texttt{extractFunctions.js}, inspect \texttt{path.parent} and \texttt{path.parentPath.parent} to determine context.

\subsection{Improvement 2: Python File Support}

\textbf{Problem:} The AST parser only handles JS/TS. Python repositories (which are common) cannot be indexed.

\textbf{Solution:} Add a Python parser in the Python AI service using the \texttt{ast} standard library module:
\begin{lstlisting}[language=Python]
import ast

def extract_python_functions(file_path: str):
    with open(file_path) as f:
        tree = ast.parse(f.read())
    
    functions = []
    for node in ast.walk(tree):
        if isinstance(node, (ast.FunctionDef, ast.AsyncFunctionDef)):
            functions.append({
                "name": node.name,
                "startLine": node.lineno,
                "endLine": node.end_lineno,
                "code": ast.get_source_segment(source, node)
            })
    return functions
\end{lstlisting}

The Node.js backend can detect \texttt{.py} files and forward them to a new \texttt{/parse-python} endpoint in the Python service, which returns the extracted functions in the same format as the JS extractor.

\subsection{Improvement 3: Call Graph / Dependency Tracking}

\textbf{Problem:} Chunks are isolated --- the system doesn't know that \texttt{isAuth.js} middleware is called before \texttt{productController.js}.

\textbf{Solution:} During AST traversal, record:
\begin{enumerate}
    \item Which functions each file imports from other files
    \item Which functions each function calls internally
    \item Which routes use which middleware (from \texttt{router.use()} calls)
\end{enumerate}

Store this as a call graph in MongoDB. When a query retrieves chunk A, also fetch chunks that A imports or that import A (one hop expansion). This gives the LLM connected context rather than isolated fragments.

\subsection{Improvement 4: BM25 Index Caching}

\textbf{Problem:} Currently, every query fetches all ChromaDB documents and rebuilds the BM25 index from scratch. For a 3000-chunk repository, this adds \textasciitilde200\,ms overhead per query.

\textbf{Solution:} Cache the BM25 index in memory keyed by repo name. Only rebuild when a new version of the repo is indexed:
\begin{lstlisting}[language=Python]
bm25_cache = {}  # module-level cache

def get_bm25_index(repo_name, documents, ids):
    if repo_name not in bm25_cache:
        tokenized = [_tokenize(doc) for doc in documents]
        bm25_cache[repo_name] = {
            "index": BM25Okapi(tokenized),
            "ids": ids
        }
    return bm25_cache[repo_name]
\end{lstlisting}

\subsection{Improvement 5: Server-Side Session History \textbf{[UPDATED]}}

\textbf{Status:} Client-side persistence is now implemented. The chat, selected repository, repository metadata, per-answer execution flows, citations, and File Viewer state are saved to \texttt{localStorage} under a single versioned key and restored automatically on refresh --- no ``Restore'' button. Corrupted or out-of-version data is cleared gracefully.

\textbf{Remaining work:} LocalStorage is per-browser. Server-side history would enable cross-device access and shareable conversations --- storing sessions keyed by \texttt{repoName} in a database and restoring them on load. The versioned storage schema was designed to make this migration straightforward.

\subsection{Improvement 6: Deployment}

\textbf{Target architecture:}
\begin{itemize}
    \item \textbf{Frontend:} Vercel (free tier, automatic CI/CD from GitHub)
    \item \textbf{Node.js Backend:} Render (free tier, Docker container)
    \item \textbf{Python AI Service:} Render (free tier, Docker container with \texttt{sentence-transformers} and \texttt{chromadb} installed)
\end{itemize}

\textbf{Challenges for deployment:}
\begin{itemize}
    \item \texttt{temp\_repos/} needs persistent disk storage (Render provides this)
    \item \texttt{chroma\_store/} needs persistent disk storage (same)
    \item MiniLM model download on first boot takes time --- need to pre-bake into Docker image
    \item Free tier services sleep after inactivity --- need a keep-alive ping
\end{itemize}

\subsection{Improvement 7: Clickable Citations \textbf{[IMPLEMENTED]}}

\textbf{Status:} Implemented in v1.0. Rather than regex-matching arbitrary file-path text, each inline-code token in an answer is matched against that message's \texttt{sourceFiles} and \texttt{sourceFunctions}. Confirmed matches become clickable: a file citation opens the File Viewer; a function citation opens its file, scrolls to the function, and briefly highlights its lines. Unmatched code spans stay inert, which avoids false-positive links. Execution Flow nodes provide the same navigation from the graph.

\subsection{Improvement 8: Rate Limiting}

\textbf{Problem:} No protection against abuse. A user could analyze hundreds of repos, consuming disk space and API credits.

\textbf{Solution:} Add rate limiting middleware to the Node.js backend:
\begin{lstlisting}[language=JavaScript]
const rateLimit = require('express-rate-limit')

const analyzeLimiter = rateLimit({
    windowMs: 60 * 60 * 1000, // 1 hour
    max: 5, // max 5 repos per IP per hour
    message: { error: 'Too many requests. Try again later.' }
})

const queryLimiter = rateLimit({
    windowMs: 60 * 1000, // 1 minute
    max: 20, // max 20 questions per minute per IP
})

router.post('/analyze', analyzeLimiter, analyzeRepo)
router.post('/query', queryLimiter, queryRepo)
\end{lstlisting}

\subsection{Improvement 9: Migrate Gemini SDK to google-genai \textbf{[NEW]}}

\textbf{Problem:} The service currently uses \texttt{google-generativeai}, which Google has deprecated. A \texttt{FutureWarning} is emitted at every startup. The newer \texttt{google-genai} package offers a cleaner client API and will be the only supported option going forward.

\textbf{Solution:}
\begin{lstlisting}[language=bash]
pip uninstall google-generativeai
pip install google-genai
pip freeze > requirements.txt
\end{lstlisting}

Then update \texttt{llm.py}:
\begin{lstlisting}[language=Python]
from google import genai

client = genai.Client(api_key=os.getenv("GEMINI_API_KEY"))

response = client.models.generate_content(
    model="gemini-2.5-flash",
    contents=prompt
)
answer = response.text
\end{lstlisting}

\subsection{Improvement 10: Fix requirements.txt \textbf{[NEW]}}

\textbf{Problem:} \texttt{rank-bm25} was installed manually and is missing from \texttt{requirements.txt}. A fresh environment setup will fail.

\textbf{Solution:} After completing the Gemini SDK migration above, regenerate the file:
\begin{lstlisting}[language=bash]
cd service
source ../.venv/bin/activate
pip freeze > requirements.txt
\end{lstlisting}

Verify the output contains at minimum: \texttt{rank-bm25}, \texttt{google-genai}, \texttt{fastapi}, \texttt{uvicorn}, \texttt{sentence-transformers}, \texttt{chromadb}, \texttt{numpy}.

\newpage

% ═══════════════════════════════════════════════════════
%  SECTION 10: INTERVIEW PREPARATION
% ═══════════════════════════════════════════════════════
\section{Interview Preparation: Questions You Will Be Asked}

\subsection{System Design Questions}

\textbf{Q: Why did you separate the backend into Node.js and Python instead of using one language?}\\[3pt]
A: Node.js has excellent tooling for the web layer --- serving HTTP, routing, file system operations, git cloning. Python has far better ML ecosystem support --- sentence-transformers, ChromaDB, BM25, LangChain all have native Python libraries. Using Python for the AI service and Node.js for the web layer lets each do what it's best at. They communicate via HTTP.

\vspace{8pt}
\textbf{Q: Why not just send the entire codebase to the LLM context?}\\[3pt]
A: Two reasons. First, a codebase with 3000+ functions would exceed the context window of any LLM. Second, even if it fit, attention dilution becomes a problem --- the LLM performs worse with irrelevant context. RAG lets us send only the 5--8 most relevant chunks, giving the LLM focused context and better answers.

\vspace{8pt}
\textbf{Q: How does your hybrid search work?}\\[3pt]
A: We run two searches simultaneously. BM25 scores all chunks based on exact keyword frequency and inverse document frequency. Vector search embeds the query and finds semantically similar chunks using cosine similarity. We merge the two ranked lists using Reciprocal Rank Fusion: $\text{score}(d) = \frac{1}{60 + \text{rank}_1(d)} + \frac{1}{60 + \text{rank}_2(d)}$. Documents appearing high in both lists score highest.

\vspace{8pt}
\textbf{Q: What is the caching strategy?}\\[3pt]
A: Two levels. First, cloned repositories are stored in \texttt{temp\_repos/} and never re-cloned if already present. Second, ChromaDB collections persist to disk in \texttt{chroma\_store/}. The \texttt{is\_already\_indexed()} function checks if a collection exists before triggering re-embedding. A repository is indexed once and served forever.

\vspace{8pt}
\textbf{Q: You have two separate AI pipelines. How are they different? \textbf{[NEW]}}\\[3pt]
A: Yes, the system has two completely independent pipelines. The RAG Q\&A pipeline indexes code chunks (function bodies and file summaries) as MiniLM embeddings in ChromaDB, and uses hybrid BM25 + vector search to retrieve relevant chunks for each user question. The git evolution pipeline is entirely separate --- it uses \texttt{git show} to get commit diffs, maps line changes to function line ranges to identify which functions changed in which commits, and sends that structured history to Gemini as plain text for summarization. Git diffs are never embedded; they go directly to the LLM.

\vspace{8pt}
\textbf{Q: How is the Execution Flow graph generated and kept accurate? \textbf{[NEW]}}\\[3pt]
A: The LLM produces it directly. The query prompt asks Gemini (in JSON mode) to return a single object with both the markdown answer and an \texttt{execution\_flow} of typed nodes and edges, so the frontend never parses a graph out of prose. Code-backed nodes carry file and function names from the retrieved context; on the server I reconcile their line ranges against the indexed function metadata so clicking a node lands on the real lines. The frontend lays the graph out with Dagre and renders it with React Flow. A sanitizer removes malformed nodes/edges, and if the model returns nothing usable, a fallback flow is synthesized from the retrieved functions so the graph is never empty. The design keeps responsibilities separate: chat explains, the Execution Flow visualizes, the File Viewer shows code, and Git Evolution shows history.

\subsection{AI/ML Questions}

\textbf{Q: Why MiniLM over OpenAI embeddings?}\\[3pt]
A: MiniLM is free and runs locally. OpenAI charges per token for embeddings. For a student project indexing thousands of functions across multiple repositories, OpenAI embeddings would incur real cost. MiniLM gives 80--90\% of the quality at zero cost.

\vspace{8pt}
\textbf{Q: What is the dimensionality of your embeddings and why does it matter?}\\[3pt]
A: MiniLM produces 384-dimensional vectors. Higher dimensionality (like OpenAI's 1536) captures more nuance but requires more storage and computation. 384 dimensions is sufficient for our use case --- we're doing code retrieval, not fine-grained semantic similarity for medical or legal documents.

\vspace{8pt}
\textbf{Q: How do you prevent hallucination?}\\[3pt]
A: Three mechanisms. First, the prompt strictly says ``answer ONLY based on the provided chunks.'' Second, it says ``only state what is clearly visible.'' Third, it instructs the LLM to explicitly report what information is missing if the chunks don't fully cover the question. This makes the LLM honest about its limitations rather than filling gaps with invented information.

\vspace{8pt}
\textbf{Q: What is Reciprocal Rank Fusion and why use it over score averaging?}\\[3pt]
A: RRF merges ranked lists based on position, not score. Score averaging fails because BM25 and vector search produce scores on completely different scales --- BM25 might score 45.2 while cosine similarity scores 0.87. Normalizing them introduces arbitrary biases. RRF only uses rank position (1st, 2nd, 3rd\ldots) which is scale-independent. A document ranked 1st in both lists always wins, regardless of the absolute scores.

\subsection{System-Specific Questions}

\textbf{Q: How does your AST parser handle different function types?}\\[3pt]
A: We handle four types: FunctionDeclaration (\texttt{function foo(){}}), FunctionExpression (\texttt{const foo = function(){}}), ArrowFunctionExpression (\texttt{const foo = () => {}}), and ClassMethod (methods in ES6 classes). The Babel traverse API calls our callback for each node type.

\vspace{8pt}
\textbf{Q: What happens when someone asks about code that's not in the top 5 chunks?}\\[3pt]
A: The LLM is instructed to explicitly state what is missing. Instead of hallucinating, it says ``Based on the provided chunks, I can see X and Y. However, the chunks don't show how Z works. To understand Z, you should look at [files].'' This is more useful than a confident but wrong answer.

\vspace{8pt}
\textbf{Q: How do you handle the focus map if all functions are anonymous?}\\[3pt]
A: We display them as \texttt{filename.js:lineNumber} instead of function names. This at least tells the user where in the file to look. The proper fix (in progress) is to infer names from context --- route handlers become ``GET /path handler'', etc.

\vspace{8pt}
\textbf{Q: How does the git evolution feature work technically? \textbf{[NEW]}}\\[3pt]
A: \texttt{gitEvolution.js} iterates over the last 20 commits using \texttt{simple-git}. For each commit, it runs \texttt{git show --unified=0} which returns a diff with zero surrounding context lines, making it easy to extract exactly which line numbers changed. For each changed line, it checks whether that line falls within any extracted function's \texttt{startLine}--\texttt{endLine} range. If it does, that commit is associated with that function. The result is an \texttt{evolutionMap} where each key is \texttt{filePath::functionName} and the value is a list of commit objects. This map is sent to Gemini, which produces a natural-language summary of how the function evolved.

\newpage

% ═══════════════════════════════════════════════════════
%  SECTION 11: TECH STACK SUMMARY
% ═══════════════════════════════════════════════════════
\section{Complete Tech Stack Reference}

\vspace{6pt}
\begin{center}
\begin{tabularx}{0.96\textwidth}{@{}>{\bfseries}lll@{}}
\toprule
\textbf{Category} & \textbf{Technology} & \textbf{Purpose} \\
\midrule
\multicolumn{3}{@{}l}{\textit{\small\color{accentblue} Frontend}} \\[2pt]
Framework    & React 18           & Component-based UI \\
Build Tool   & Vite               & Dev server and bundling \\
Routing      & React Router v6    & Client-side navigation \\
Styling      & Tailwind CSS v3    & Utility-first CSS \\
HTTP Client  & axios              & API calls to backend \\
Markdown     & react-markdown     & Render LLM responses \\
Graph Viz    & React Flow + Dagre & Execution Flow rendering + layout \textbf{[NEW]} \\
Syntax HL    & Shiki              & Code highlighting \textbf{[NEW]} \\
\midrule
\multicolumn{3}{@{}l}{\textit{\small\color{accentblue} Node.js Backend}} \\[2pt]
Framework    & Express.js         & HTTP server and routing \\
Git Ops      & simple-git         & Clone repos; run git show for evolution \textbf{[UPDATED]} \\
JS/TS Parser & @babel/parser      & Parse code to AST \\
AST Walker   & @babel/traverse    & Walk AST nodes \\
HTTP Client  & axios              & Call Python service \\
Environment  & dotenv             & Load .env variables \\
Repo Cleanup & repoController.js  & Bound temp\_repos to recent clones \textbf{[NEW]} \\
CORS         & cors               & Allow cross-origin requests \\
Git History  & gitEvolution.js    & Map commit diffs to functions \textbf{[NEW]} \\
MVC          & repoController.js  & Separated business logic from routing \textbf{[NEW]} \\
\midrule
\multicolumn{3}{@{}l}{\textit{\small\color{accentblue} Python AI Service}} \\[2pt]
Framework    & FastAPI                & Python HTTP service \\
Embed Model  & MiniLM-L6-v2          & Text to vector (384-dim) \\
Embed Lib    & sentence-transformers  & Run MiniLM locally \\
Vector DB    & ChromaDB              & Store and search embeddings \\
Keyword      & rank-bm25             & BM25 keyword retrieval \\
Merge        & RRF (custom)          & Combine BM25 + vector \\
LLM          & Gemini 2.5 Flash      & Answer + execution flow + evolution summaries \textbf{[UPDATED]} \\
LLM SDK      & google-generativeai   & Python client for Gemini (migration to google-genai pending) \textbf{[UPDATED]} \\
Numerical    & numpy                 & Vector operations \\
\midrule
\multicolumn{3}{@{}l}{\textit{\small\color{accentblue} Infrastructure}} \\[2pt]
Version Ctrl & Git + GitHub & Source control \\
Runtime (BE) & Node.js      & JavaScript server runtime \\
Runtime (AI) & Python 3.11  & AI service runtime \\
\bottomrule
\end{tabularx}
\end{center}

\vspace{1.2cm}

\begin{center}
\textcolor{medgray}{\rule{0.5\linewidth}{0.4pt}}
\end{center}

\vspace{0.4cm}

\noindent\textbf{Repository:}\quad\href{https://github.com/Sahil0282/repo-explorer}{github.com/Sahil0282/repo-explorer}

\vspace{0.4cm}

\noindent\textbf{Author:}\quad Sahil Arun Dhawane\quad|\quad Vishwakarma Institute of Technology, Pune\\[0.2cm]
\noindent\textbf{Program:}\quad B.Tech Electronics \& Telecommunications \textit{(Expected 2027)}

\end{document}